% FILE: 01_research_question.tex
% Project: standards
% Thesis Role: public-standards-compliance-and-gravity-field-mapping

\section{Research Question}
\label{sec:standards-rq}

\subsection{Problem Statement}
\label{sec:standards-rq-problem}

Process intelligence systems routinely issue claims about operational conformance, audit
fitness, and process soundness. In practice, these claims are made in proprietary formats
against proprietary models, with no named mathematical referent and no mechanism by which
an independent party---a buyer's advisor in an M\&A data room, a regulator, a competing
vendor---can replicate the verification. The result is an ecosystem where process evidence
is structurally indistinguishable from marketing assertion.

The problem this project addresses is: \emph{what formal apparatus is required to convert
a process evidence claim into a board-admissible, cryptographically verifiable fact?} The
answer demanded by the research program is a complete mapping from each claim type to a
named public standard, a named mathematical formulation within that standard, a named
validation threshold, and a BLAKE3 receipt chain that binds the claim to its evidence.
Without this mapping, the Evidence lattice \texttt{Evidence<T, State, Witness>} has no
external referent; it is a type system that witnesses only itself.

\subsection{Old-Regime Assumption Challenged}
\label{sec:standards-rq-old-regime}

The old-regime assumption is that standards compliance is a documentation exercise: a
vendor asserts conformance to XES, OCEL, or BPMN by stating it in a specification sheet.
No formal mapping is required; no mathematical formulation is quoted; no validation
threshold is named. The standard is treated as a badge rather than a constraint.

This assumption has three structural defects. First, it permits process log laundering:
event data can be formatted as XES while violating chronological invariance
($\forall i,\ t(e_{i+1}) \ge t(e_i)$), with no enforcement mechanism to detect the
violation. Second, it permits phantom object injection in OCEL logs---events bound to
object identifiers that do not appear in the object set---because the bipartite graph
invariants of $\mathcal{L}=(\mathcal{E}, \mathcal{O}, \text{E2O}, \text{O2O}, \Delta)$
are never checked. Third, it permits WF-net soundness claims without executing the
Karp-Miller Coverability Graph construction that would detect dead transitions or
non-terminating markings.

The adversarial test report \texttt{ocel-adversarial-gravity.md} (identifier
OCEL-XES-ADV-V30.1.1) documents each of these vulnerabilities in XES and OCEL native
formats under the old regime.

\subsection{Hypothesis}
\label{sec:standards-rq-hypothesis}

\textbf{Hypothesis:} If each public standard is formally mapped to (a) a named
mathematical formulation, (b) a Blue River Dam validation threshold, (c) a type-law
surface in the \texttt{Evidence<T, State, Witness>} lattice, and (d) a BLAKE3 receipt
schema, then process evidence claims become independently replicable, machine-verifiable
facts rather than vendor assertions---and this replicability is both a customer protection
and a competitive weapon.

The hypothesis further claims that this mapping, when applied to 16 standards, creates a
gravitational field: the foundry's architecture is held in place by the mutual
gravitational pull of independently verifiable specifications, such that no single
proprietary element is load-bearing. This is the architectural basis for the reverse
lock-in doctrine.

\subsection{Success Criteria}
\label{sec:standards-rq-success}

The graduation gate for this research project, as recorded in checkpoint
\texttt{chk-public-standards-crosswalk-001}, specifies five criteria:

\begin{enumerate}
  \item \textbf{Type-safety mapping}: All 16 standards are mapped to the
        \texttt{Evidence<T, State, Witness>} lattice, as verified against the Type-Law Atlas
        at \texttt{sources/wasm4pm-compat/type-law-atlas.md}.
  \item \textbf{PM4Py Oracle Map alignment}: Mappings are aligned with the limits and
        assumptions identified in the PM4Py Oracle Map at
        \texttt{sources/pm4py/oracle-map.md}.
  \item \textbf{Audit completeness}: The audit matrix \texttt{audit\_\_standards\_coverage.md}
        registers green (VERIFIED) status for all 16 standards.
  \item \textbf{No template files}: All placement files contain fully realized schemas,
        mathematical formulations, and execution rules.
  \item \textbf{Link integrity}: All cross-references are written as absolute markdown
        links with no broken paths.
\end{enumerate}

All five criteria were satisfied at checkpoint timestamp \texttt{2026-05-31T23:59:59Z},
graduating the standards crosswalk to ALIVE status.
